% %%%

%%%   Самарский государственный технический университет

%%%   Кафедра прикладной математики и информатики

%%%

%%%   Преамбула для статьи в журнал "Вестник Самарского государственного

%%%   технического университета. Серия: «Физико-математические науки»"

%%%   Ориентирована под MikTEx 2.4 for Win32

%%%

%%%   Хотя сам журнал собирается под GNU/Linux на дистрибутиве TeXLive



\documentclass[11pt,a4paper,twoside]{article}



\usepackage[cp1251]{inputenc}

\usepackage[T2A]{fontenc}

\usepackage[english,russian]{babel}

\usepackage{samgtubib}

\usepackage[dvips]{graphicx}

\usepackage[multiple]{footmisc}



\usepackage{samgtu}

\renewcommand{\VestnikYear}{2009} % Год издания Вестника

\renewcommand{\VestnikNumber}{2\,(19)} % Номер Вестника

\renewcommand{\VestnikMonth}{Ноябрь} % Месяц выхода Вестника

\renewcommand{\VestnikData}{01.11.09} % Дата подписания Вестника в печать

\renewcommand{\VestnikUPL}{26,64} % Усл. печ. л.

\renewcommand{\VestnikUIL}{25,73} % Уч.-изд. л.

\renewcommand{\VestnikRegNumber}{479/09} % Регистрационный номер

\renewcommand{\VestnikOrder}{994} % Номер заказа







%

% \begin{document}

% \tableofcontents

% \newpage



\renewcommand{\firstpage}{6}

\renewcommand{\lastpage}{17}

%\Partition{Дифференциальные уравнения}{} % Раздел Вестника, в который подаётся статья

% Названия разделов см. на сайте при подаче заявки





%%%%%%%%%%%%%%%%%%%%%%%%%%%%%%%%%%%%%%%%%%%%%%%%%%%%%%%%%%%%%%%%%%%%%%%%%%%%%%%%%%%%

%Информация о статье и об авторах на русском и английском языках



% Номер УДК

\udk{517.957}%Обработка текста. Подготовка текстов : Языки программирования. Компьютерные языки

\msc{35K65}% Коды MSC см. http://www.ams.org/msc/



\title{О прямой и обратной задачах для уравнений Хоффа на графе}% Полный заголовок

{О прямой и обратной задачах \dots}% Заголовок, который идёт в верхний колонтитул



\author{Г.\,А.~Свиридюк, А.\,А.~Баязитова}% Автор(ы) для заголовка, если авторов несколько укать \inst

{С\,в\,и\,р\,и\,д\,ю\,к~Г.\,А.,\ Б\,а\,я\,з\,и\,т\,о\,в\,а~А.\,А.}% Автор(ы) для оглавления и колонтитула через \,



\institute{Южно-Уральский государственный университет, \\  454080 Челябинск, пр. Ленина, 76.}

% Если несколько, то так  \institute{ Организация 1 \and Организация 2  \and Организация 3}



\email{E-mails}{ridyu@math.susu.ac.ru, alfiya@74.ru} %E-mail's авторов



\otherinf{Свиридюк Георгий Анатольевич "--- зав. кафедрой уравнений математической физики; д.ф.-м.н., профессор.\\

Баязитова Альфия Адыгамовна "--- аспирант кафедры уравнений математической физики.}



\keywords{уравнения Хоффа на графе, фазовое пространство, обратная задача}



\workabstract{% Здесь должна быть краткая аннотация статьи, например такая:

Уравнения Хоффа, определенные на графе, описывают динамику выпучивания конструкции из двутавровых балок. Получено обобщение прямой задачи "--- задачи Коши. Впервые рассмотрена обратная коэффициентная задача, моделирующая эксперимент, в результате которого при дополнительных измерениях изучается не только динамика выпучивания конструкции, но и свойства материалов балок. Показано существование единственного решения этой задачи.}



\engtitle{On Direct and Inverse Problems for the Hoff Equations on Graph}



\engauthor{G.\,A.~Sviridyuk, A.\,A.~Bayazitova}{S\,v\,i\,r\,i\,d\,y\,u\,k~G.\,A., B\,a\,y\,a\,z\,i\,t\,o\,v\,a~A.\,A.}





\enginstitute{South Ural State University,\\

76, prosp. Lenina,  Chelyabinsk, 454080.}% Место работы каждого из авторов на англ. языке

% Если несколько, то так  \enginstitute{ Организация 1 \and Организация 2  \and Организация 3}



\otherenginfm{Sviridyuk Georgy Anatolevich, Dr. Sci. $($Phys. \& Math.$)$, Prof.,  Head of Dept. of Mathema\-tical Physics Equations.\\

Bayazitova Alfiya Adygamovna, Postgraduate Student, Dept. of Mathematical Physics Equations.}



\engabstract{The Hoff equations defined on graph describe the dynamics of H-beams construction buckling. Generalization of the direct problem which is the Cauchy problem is obtained. For the first time the inverse coefficient problem is studied which is modeling the experiment that allows with additional measurements not only to study the construction buckling dynamics but also the characteristics of beam material. The unique solution of this problem is demonstrated.}



\engkeywords{the Hoff equations on graph, phase space, an inverse problem}



\date{19/I/2009} % Дата поступления статьи в редакцию.

\revisiondate{19/II/2009} % В окончательном виде



%%%%%%%%%%%%%%%%%%%%%%%%%%%%%%%%%%%%%%%%%%%%%%%%%%%%%%%%%%%%%%%%%%%%%%%%%%%%%%%%%%%%





\maketitle





\Section[n]{Введение}

 Пусть $\mathbf{G}=\mathbf{G}(\mathfrak{V};\mathfrak{E})$ "--- конечный связный ориентированный \linebreak граф,

 где $\mathfrak{V}=\{V_i\}$ "--- множество вершин, а $\mathfrak{E}=\{E_j\}$ "--- множество рёбер,

 причём каждому ребру $E_j$ поставлены в соответствие два положительных числа $l_j$, $d_j\in \mathbb{R}_{+}$,   которые в контексте нашей задачи будут иметь физический смысл длины и площади поперечного сечения ребра соответственно.

 На каждом ребре $E_j$ задается урав\-не\-ние Хоффа

 \begin{equation}\label{svirid:0.1}

 \lambda_ju_{jt}+u_{jxxt}=\alpha_ju_j+\beta_ju_j^3,

 \end{equation}

моделирующее выпучивание двутавровой балки, где параметр

$\lambda_j\in \mathbb{R}_+$ со\-от\-вет\-ству\-ет нагрузке на балку, а

параметры $\alpha_j$, $\beta_j\in \mathbb{R}$ характеризуют

свойства материала; переменные $x\in(0, l_j)$, $t\in

\mathbb{R}$.



Под прямой задачей понимается поиск вектор"=функции

$u=(u_1,\, u_2,\, \dots,$ $u_j,\, \dots)$, каждая компонента которой

$u_j=u_j(x,t)$

удовлетворяет уравнению (\ref{svirid:0.1}) на ребре $E_j$, а в вершинах

$\mathfrak{V}$ компоненты удовлетворяют условию

непрерывности

\begin{equation}\label{svirid:0.2}

\begin{split}

u_j(0,\, t)=u_k(0,\, t)=u_m(l_m,\, t)=u_n(l_n,\, t),\\

E_j,\, E_k\in \mathfrak{E}^\alpha(V_i), \quad E_m,\, E_n\in

\mathfrak{E}^\omega(V_i)

\end{split}

\end{equation}

и условию <<баланса потока>>

\begin{equation}\label{svirid:0.3}

\sum\limits_{E_j\in \mathfrak{E}^\alpha(V_i)}d_ju_{jx}(0,

t)-\sum\limits_{E_k\in \mathfrak{E}^\omega(V_i)}d_ku_{kx}(l_k,\, t)=0,

\end{equation}

где через $\mathfrak{E}^{\alpha(\omega)}(V_i)$ обозначено

множество ребер, выходящих из вершины $V_i$ (вхо\-дя\-щих в вершину

$V_i$). Кроме того, искомые компоненты должны удовлетворять

начальным условиям Коши

\begin{equation}\label{svirid:0.4}

u_j(x,\, 0)=u_{j0}(x), \quad x\in(0,\, l_j).

\end{equation}

При обратной задаче неизвестными являются не только компоненты

$u_j=u_j(x, t)$, но и коэффициенты $\alpha_j$ и $\beta_j$. Для

решения обратной задачи не\-об\-хо\-ди\-мо вводить дополнительные условия,

о них речь пойдёт далее.



Условие (\ref{svirid:0.2}) требует, чтобы вектор-функция $u=u(x,\, t)$ была

не\-пре\-рыв\-ной в вершинах графа. Отметим, что в~контексте

этого условия вы\-ра\-же\-ния <<отсутствовать>>  и <<быть равным нулю>>

имеют различный смысл. Например, если в вершину $V_i$ все ребра

входят, то левые два равенства в (\ref{svirid:0.2}) именно <<отсутствуют>> , а не

<<равны нулю>>. Условия (\ref{svirid:0.3}) "--- аналог условий Кирхгоффа "---

превращаются в условия Неймана, если граф $\mathbf{G}$ состоит из

единственного ребра и двух вершин, причем условия (\ref{svirid:0.2}) в данном

случае <<отсутствуют>>. Если же граф $\mathbf{G}$ состоит из

единственного ребра и единственной вершины, то условия (\ref{svirid:0.2}),

(\ref{svirid:0.3}) превращаются в условия согласования.



Прямая задача представляет собой модель для изучения поведения

на\-гру\-жен\-ной конструкции из двутавровых балок, а обратная задача

мо\-де\-ли\-ру\-ет экс\-пе\-ри\-мент, в результате которого при

дополнительных из\-ме\-ре\-ни\-ях изучается не только динамика

конструкции, но и свойства ма\-те\-риа\-лов балок. Обе модели

представлены дифференциальными уравнениями с

част\-ны\-ми производными, за\-дан\-ны\-ми на графе. Дифференциальные урав\-не\-ния на графах "--- сравнительно

новая область математики, возникшая в конце прошлого века.

Первая мо\-но\-гра\-фия по этой проблематике вышла в 2004~г.~[1].

Обратные задачи для не\-клас\-си\-чес\-ких уравнений математической физики рассмотрены в~[2].

Уравнение Хоффа от\-но\-сит\-ся к уравнениям соболевского

типа, ис\-сле\-до\-ва\-ния которых в настоящее время пе\-ре\-жи\-ва\-ют

пору бурного рас\-цве\-та (см. исторический обзор в~[3]). Урав\-не\-ния

со\-бо\-лев\-ско\-го типа на графах впервые рассмотрены в [4]. В

[5] впервые поставлена и изучена обратная задача для линейных

уравнений соболевского типа. В [6] впервые исследована

обратная задача для уравнений Баренблатта"--~Желтова"--~Кочиной на графе.

Наконец, в [7] изучена задача (\ref{svirid:0.2})--(\ref{svirid:0.4})

для урав\-не\-ний (\ref{svirid:0.1}) при более простых предположениях на

данные задачи, именно $d_j=1$, $\alpha_j=\alpha$, $\beta_j=\beta$,

$\lambda_j=\lambda$. Обратная задача рассмативается впервые.



\Section{Прямая задача}

%\numberwithin{equation}{section}

Следуя [6], введём в рассмотрение

гильбертово пространство $\mathbf{L_2(G)}=\{g=(g_1,\, g_2,\, \dots,\, g_j,\, \dots)$, $g_j\in L_2(0,\, l_j)\}$ со скалярным произведением

$$\langle g,\, h\rangle=\sum\limits_jd_j\int_0^{l_j}g_jh_jdx$$ и~банахово пространство

$\mathbf{\mathfrak{U}}=\{u=(u_1, u_2, \dots, u_j, \dots): u_j\in

W^1_2(0, l_j)$ и вы\-пол\-не\-но (\ref{svirid:0.2}) с нормой

$$\|u\|_{\mathbf{\mathfrak{U}}}^2=\sum\limits_jd_j\int_0^{l_j}(u^2_{jx}+u_j^2)dx.$$

Отметим, что условие (\ref{svirid:0.2}) имеет смысл в силу абсолютной

непрерывности ком\-по\-нен\-тов $u_j$, а пространство

$\mathbf{\mathfrak{U}}$ плотно и компактно вложено в

$\mathbf{L_2(G)}$. Обозначим через $\mathbb{\mathfrak{F}}$

сопряженное к $\mathbf{\mathfrak{U}}$ относительно $\langle\cdot,

\cdot\rangle$ банахово про\-стран\-ство. Очевидна плотность и компактность

вложения $\mathbf{\mathfrak{U}\hookrightarrow\mathfrak{F}}$.



Далее, выберем произвольные константы $a_j\in\mathbb{R}_+$ и

построим оператор $A:\mathfrak{U}\rightarrow\mathfrak{F}$:

$$

\langle Au, v\rangle=\sum\limits_jd_j\int\limits_0^{l_j}(u_{jx}v_{jx}+a_ju_jv_j)dx,

\quad u,\, v\in\mathbf{\mathfrak{U}}.

$$

Как нетрудно заметить,

оператор $A$ самосопряжён, положительно оп\-ре\-де\-лён и является

топлинейным изоморфизмом пространств $\mathbf{\mathfrak{U}}$ и

$\mathbf{\mathfrak{F}}$. Ввиду компактности вложения

$\mathbf{\mathfrak{U}\hookrightarrow\mathfrak{F}}$ спектр

$\sigma(A)$ оператора $A$ положителен, дискретен,

ко\-неч\-но\-кра\-тен и сгущается только к $+\infty$.



Теперь построим операторы $L,\, M:\mathfrak{U}\rightarrow\mathfrak{F}$:

\begin{gather*}

\langle Lu, v\rangle=\sum\limits_jd_j(\lambda_j+a_j)\int_0^{l_j}u_jv_jdx-\langle Au, v\rangle,

\\

\langle Mu, v\rangle=\sum\limits_j\alpha_jd_j\int_0^{l_j}u_jv_jdx.

\end{gather*}

Очевидно, операторы $L,\, M\in\mathbf{\mathfrak{L(U;\,F)}}$ (т.\,е.

линейны и~непрерывны), при\-чём опе\-ра\-тор $L$ фредгольмов (т.\,е.

${\rm ind} \, L=0$), а~оператор $M$ компактен. Напомним (см.

[7, гл.~4]), что оператор $M$ называется $(L,\, 0)$-ограниченным,

если он $(L,\, \sigma)$-ограничен  и~точка

$\infty$ является устранимой особой точкой \mbox{$L$-резольвенты}

оператора $M$.

\begin{lemma}[1.1]

Оператор $M$ $(L,0)$-ограничен, если

\begin{itemize}

\item[\rm (i)] $\ker L=0;$



\item[\rm (ii)]  $\ker L\neq \{0\},$ $\alpha_j\neq 0$ при любом $j$ и все

$\alpha_j$ имеют одинаковый знак.

\end{itemize}

\end{lemma}



\begin{proof}

Утверждение (i) очевидно, т.\,к. в этом случае су\-щест\-ву\-ет

оператор $L^{-1}\in\mathbf{\mathfrak{L(U;\,F)}}$ (детали см. в~замечании 4.1.1 [8]).



Пусть $\ker L\neq \{0\}$ и $\alpha_j\in\mathbb{R}_+$. Тогда

билинейная форма

$$

[h,\,g]=\sum\limits_j\alpha_jd_j\int_0^{l_j}h_jg_jdx

$$

задаёт в $\mathbf{L_2(G)}$ скалярное произведение, эквивалентное

$\langle \cdot, \cdot\rangle$. В силу фред\-голь\-мо\-вос\-ти оператора $L$

${\rm codim} \, {\rm im} \ L={\rm dim} \, \ker L$, причём,

как нетрудно заметить, ядро $\ker L$ ортогонально образу

${\rm im \, }L$ относительно $\langle \cdot, \cdot\rangle$. Возьмём вектор

$\psi\in\ker L\backslash \{0\}$ и рассмотрим

$$

\langle M\psi, \psi\rangle=\sum\limits_j\alpha_jd_j\int_0^{l_j}\psi_j^2dx=[\psi, \psi]>0.

$$

Отсюда следует, что $M\psi\neq {\rm im \,}L$ при любом векторе

$\psi\in\ker L\backslash\{0\}$. Это в свою очередь означает, что

ни один вектор $\psi\in\ker L\backslash\{0\}$ не имеет

\mbox{$M$-присоединён}\-ных векторов, что в~силу теоремы 4.6.1.~[8]

завершает до\-ка\-за\-тель\-ство леммы в случае $\alpha_j\in\mathbb{R}_+$.

Если $\alpha_j\in\mathbb{R}_-$, доказательство аналогично.~\end{proof}



Теперь построим оператор

$$\langle N(u), v\rangle=\sum\limits_j\beta_jd_j\int_0^{l_j}u_j^3v_jdx$$

и убедимся, что он действует из пространства $\mathfrak{U}$ в

пространство $\mathfrak{F}$. Для этого построим вспомогательное

пространство $\mathbf{L_4(G)}=\{g=(g_1, g_2, \dots,$ $g_j, \dots):\,g_j\in

L_4(0,l_j)\}$. Очевидно, имеют место плотные и непрерывные

вложения

$\mathfrak{U}\hookrightarrow\mathbf{L_4(G)}$ $\hookrightarrow\mathbf{L_2(G)}$.

Обозначим через $\mathbf{L_4^*(G)}$ сопряженное к

$\mathbf{L_4(G)}$ относительно двой\-ствен\-нос\-ти $\langle\cdot, \cdot\rangle$

пространство. Пространство $\mathbf{L_4^*(G)}$ топ\-ли\-ней\-но

изоморфно про\-стран\-ству

$\mathbf{L_{\frac{4}{3}}(G)}=\{g=(g_1, g_2, \dots, g_j,$ $\dots):\,g_j\in

L_{\frac{4}{3}}(0, l_j)\}$. Норма в пространствах $\mathbf{L_p(G)}$

задаётся следующим образом:

$$\|u\|_p=\sum\limits_jd_j\biggl(\int_0^{l_j}|u_j|^pdx\biggr)^{\frac{1}{p}}.$$

Поэтому, в силу неравенства Гёльдера, получим

$$|\langle N(u), v\rangle|\leq \max\limits_j \{|\beta_j|\} \ \|u\|_4^3 \ \|v\|_4,$$

т.\,е. действие оператора

$N:\mathbf{L_4(G)}\rightarrow\mathbf{L_{\frac{4}{3}}(G)}$ имеет

место.

Действие опе\-ра\-то\-ра $N:\mathfrak{U}\rightarrow\mathfrak{F}$

имеет место в силу вложения

$\mathfrak{U}\hookrightarrow\mathbf{L_4(G)}$, из которого вытекает

вложение $\mathbf{L_{\frac{4}{3}}(G)}\hookrightarrow\mathfrak{F}$.

\begin{lemma}[1.2]

При любых $\beta_j\in\mathbb{R}$ оператор

$N\in\textsl{C}^\infty(\mathfrak{U;F})$.

\end{lemma}

\begin{proof}

Фиксируем точку $u\in\mathfrak{U}$ и рассмотрим про\-из\-вод\-ную

Фре\-ше $N'_u$ оператора $N$ в точке $u$:

$$\langle N'_u (v), w\rangle=3\sum\limits_j\beta_jd_j\int_0^{l_j}u_j^2v_jw_jdx.$$

Отсюда аналогично предыдущему имеем

$$|\langle N'_u(v), w\rangle|\leq 3 \max\limits_j \{|\beta_j|\} \ \|u\|_4^2 \ \|v\|_4 \ \|w\|_4,$$

т.\,е. $N'_u\in\mathfrak{L(U;\,F)}$ при фиксированном $u$.

Непрерывность второй и треть\-ей про\-из\-вод\-ных Фреше доказывается

аналогично, остальные производные рав\-ны нулю.

\end{proof}



Итак, мы редуцировали задачу (\ref{svirid:0.1})--(\ref{svirid:0.4}) к задаче Коши

\begin{equation}\label{svirid:1.1}

u(0)=u_0

\end{equation}

для полулинейного уравнения соболевского типа

\begin{equation}\label{svirid:1.2}

L\dot{u}=Mu+N(u),

\end{equation}

где операторы $L,\,M\in\mathfrak{L(U; F)}$,

$N\in\textsl{C}^\infty\mathfrak{(U;\,F)}$, причём оператор $L$

фред\-голь\-мов, а оператор $M$ $(L,\,0)$-ограничен. Вектор"=функцию

$u\in\textsl{C}^1((-\tau, \tau);\, \mathfrak{U})$, удов\-лет\-во\-ряю\-щую

уравнению (\ref{svirid:1.2}) при некотором $\tau\in\mathbb{R}_+$, назовём

ре\-ше\-ни\-ем этого уравнения, а если решение вдобавок удовлетворяет

условию (\ref{svirid:1.1}), то будем называть его решением задачи (\ref{svirid:1.1}), (\ref{svirid:1.2}).



Если $\ker L=\{0\}$, то уравнение (\ref{svirid:1.2}) тривиально редуцируется к

эк\-ви\-ва\-лент\-но\-му ему уравнению

\begin{equation}\label{svirid:1.3}

\dot{u}=F(u),

\end{equation}

где оператор $F=L^{-1}(M+N)\in\textsl{C}^\infty(\mathfrak{U})$ по

построению. Существование единст\-вен\-но\-го локального решения

$u\in\textsl{C}^1((-\tau, \tau);\,\mathfrak{U})$ задачи (\ref{svirid:1.1}), (\ref{svirid:1.2})

при любом $u_0\in\mathfrak{U}$ "--- результат классической теоремы

Коши (см., например, [9, гл.~IV, $\S 1$]). Другое дело, если $\ker

L\neq \{0\}$. В этом случае полезным оказывается следующее

понятие.

\begin{definition}[1.1]

Множество $\mathfrak{P}\subset\mathfrak{U}$

называется {\it фазовым пространством урав\-не\-ния} (\ref{svirid:1.2}), если:



\begin{itemize}

    \item [(i)] любое решение $u=u(t)$ уравнения (\ref{svirid:1.2}) лежит в $\mathfrak{P}$

как траектория, т.\,е. $u(t)\in\mathfrak{P}$ $t\in(\tau, \tau)$;

    \item [(ii)] при любом $u_0\in\mathfrak{P}$ существует единственное

решение задачи (\ref{svirid:1.1}), (\ref{svirid:1.2}).

\end{itemize}

\end{definition}



Поскольку решение задачи (\ref{svirid:1.1}), (\ref{svirid:1.3}) в случае $\ker L=\{0\}$

яв\-ля\-ет\-ся решением задачи (\ref{svirid:1.1}), (\ref{svirid:1.2}) (и~наоборот), то фазовым

пространством урав\-не\-ния (\ref{svirid:1.2}) в~данном случае служит все

пространство $\mathfrak{U}$. Рассмотрим слу\-чай $\ker L\neq\{0\}$.

В~[6] для изучения этого случая была использована теорема о

расщеплении (теорема 4.1.1. [8]). Однако для наших целей более

подходит другой метод, развитый в [10, 11] и примененный в~[12].

Со\-глас\-но ему отождествим ядро и коядро оператора $L$, т.\,е. положим

${\rm coker}\, L=\ker L$, а через ${\rm coim \, } L$ обозначим

ортогональное в смысле $\mathbf{L_2(G)}$ дополнение к~ядру $\ker

L$ в пространстве $\mathfrak{U}$, т.\,е. положим ${\rm coim \, }

L=(\ker L)^\bot$. По по\-строе\-нию ${\rm coim \, } L$ плотно и~непрерывно

вложено в~${\rm im \, } L$. Обозначим через $P(Q)$

про\-ек\-тор на ${\rm coim \, } L$ $({\rm im \, } L)$ вдоль $\ker L$

$({\rm coker \, } L)$, а через $L_1$ сужение оператора $L$ на

${\rm coim \, } L$. В силу теоремы Банаха существует оператор

$L_1^{-1}\in\mathfrak{L}({\rm im \, } L; {\rm coim \, } L)$, поэтому

уравнение (\ref{svirid:1.2}) можно редуцировать к эк\-ви\-ва\-лент\-ной ему

системе уравнений

\begin{equation}\label{svirid:1.4}

0=(\mathbb{I}-Q)(M+N)(u^0+u^1), \dot{u}^1=L^{-1}Q(M+N)(u^0+u^1),

\end{equation}

где $u^1=Pu$, $u^0=u-u^1$. Заметим, что в [7] уравнение (\ref{svirid:1.2}) тоже

ре\-ду\-ци\-ру\-ет\-ся к (\ref{svirid:1.4}), однако проекторы $P$ и $Q$ строятся по

операторам $L$ и $M$ однозначно. Тем не менее в [8, п.\,4.5]

показано, что в случае $(L,\,0)$-ограниченности оператора $M$ оба

подхода приводят к одинаковым ре\-зуль\-та\-там.



Рассмотрим множество

$\mathfrak{M}=\{u\in\mathfrak{U}:(\mathbb{I}-Q)(Mu+N(u))=0\}$.

Поскольку любое решение уравнения (\ref{svirid:1.2}) с необходимостью лежит во

множестве $\mathfrak{M}$ как траектория, то это множество

является очевидным кан\-ди\-да\-том на роль фазового пространства

уравнения (\ref{svirid:1.2}).

Приведём дос\-та\-точ\-ные условия того, что

$\mathfrak{M}$ яв\-ля\-ет\-ся фазовым про\-стран\-ством. Пусть

$\mathfrak{M}\neq \emptyset$, т.\,е. существует точка

$u_0\in\mathfrak{M}$, причём оператор

$(\mathbb{I}-Q)(M+N'_{u_0})(\mathbb{I}-P):\ker

L\rightarrow{\rm coker \ } L$ является топлинейным изоморфизмом.

Тогда в силу теоремы о неявной функции существуют ок\-рест\-нос\-ти

$\mathcal{O}\subset\ker L$, $\mathcal{O}^1\subset{\rm coim \, } L$

точек $u_0^0=(\mathbb{I}-P)u_0$, $u_0^1=Pu_0$ соответственное и

оператор

$\mathcal{B}\in\textsl{C}^\infty(\mathcal{O}^1;\,\mathcal{O}^0)$

такой, что $u_0^0=\mathcal{B}(u_0^1)$.

Построим оператор

$\mathcal{D}=\mathbb{I}+\mathcal{B}:\mathcal{O}^1\rightarrow\mathfrak{M}$,

$\mathcal{D}(u_0^1)=u_0$, который вкупе с множеством

$\mathcal{O}^1$ задаёт карту множества $\mathfrak{M}$, причём

$\mathcal{D}^{-1}$ равен сужению проектора $P$ на

$\mathcal{D}[\mathcal{O}^1]=\mathcal{O}\subset\mathfrak{M}$.

По построению оператор

$\mathcal{D}\in\textsl{C}^\infty(\mathcal{O}^1;\,\mathcal{O})$.

Подействуем производной Фреше $\mathcal{D}'_{u^1}$ на нижнее

уравнение системы (\ref{svirid:1.4}), получим уравнение (\ref{svirid:1.3}), определённое на

$\mathcal{O}$, где $F=\mathcal{D}'_{u_1}L_1^{-1}Q(M+N)(u)$,

$u=\mathcal{D}(u^1)$.

Однозначная локальная разрешимость задачи

(\ref{svirid:1.1}), (\ref{svirid:1.3}) и в данной ситуации "--- ре\-зуль\-тат все той же

классической теоремы Коши [9, гл. IV, $\S 1$].



Итак, поскольку найденное решение $u=u(t)$ задачи (\ref{svirid:1.1}), (\ref{svirid:1.3})

яв\-ля\-ет\-ся решением задачи (\ref{svirid:1.1}), (\ref{svirid:1.2}), то достаточные условия

сводятся к следующим. Во-первых, $\mathfrak{M}\neq\emptyset$, а

во-вторых, в точке $u_0\in\mathfrak{M}$ существует карта $(D,

\mathcal{O}^1)$ такая, что $\mathcal{D}^{-1}$ есть сужение

проектора $P$ на $\mathcal{O}=\mathcal{D}[\mathcal{O}^1]$.



Немного отходя от стандарта (см., например, [9, гл. II]) назовём

$\mathfrak{M}$ банаховым $\textsl{C}^\infty$-многообразием в точке

$u_0$, моделируемым пространством ${\rm coim \, } L$, если

вы\-пол\-не\-ны оба условия. Если множество $\mathfrak{M}$ является

ба\-на\-хо\-вым $\textsl{C}^\infty$-многообразием, моделируемым

пространством ${\rm coim \, } L$ в каждой своей точке, то мы

назовём его  банаховым

$\textsl{C}^\infty$-мно\-го\-об\-ра\-зи\-ем.

И наконец, банахово

$\textsl{C}^\infty$-многообразие $\mathfrak{M}$ назовём простым

многообразием, если су\-же\-ние проектора $P$ на $\mathfrak{M}$

является биекцией.

\begin{lemma}[1.3]

Пусть выполнено условие {\rm (ii)} леммы {\rm 1.1} и все

ненулевые $\beta_j$ имеют тот же знак{\rm,} что и $\alpha_j.$ Тогда

множество $\mathfrak{M}$ "--- простое многообразие.

\end{lemma}



\begin{proof}

Выберем в ядре $\ker L$ ортонормированный (в смысле

$\langle\cdot, \cdot\rangle$) базис, т.\,е. $\ker L={\rm span \, } \{\chi_k:

k=1, 2, \dots, l\}$, и отождествим его с базисом в ${\rm coker \, } L$.

Тогда множество $\mathfrak{M}$ может быть записано в виде

$\mathfrak{M}=\{u\in\mathfrak{U}:\langle Mu+N(u), \chi_k\rangle=0,

k=1, 2, \dots, l\}$. Фиксируем точку $u^1\in{\rm coim \, } L$ и

рассмотрим оператор $S:\ker L\rightarrow\ker L$, задаваемый

формулой

$$S(u^0)=\sum\limits_{k=1}^l\langle (M+N)(u^1+u^0), \chi_k\rangle\chi_k.$$

По построению оператор $S\in\textsl{C}^\infty(\ker L; \ker L)$.



Пусть $\alpha_j\in\mathbb{R}_+$. Покажем, что оператор $S$ строго

монотонен, то есть $\langle S(u^0)-S(v^0), u^0-v^0\rangle>0$, если $u^0\neq v^0$.

Действительно,

\begin{multline*}

 \langle S(u^0)-S(v^0), u^0-v^0\rangle= \\

=\langle(M+N)(u^1+u^0)-(M+N)(u^1+v^0), (u^0+u^1)-(v^0+u^1)\rangle=\\

=\sum\limits_j\alpha_jd_j\int_0^{l_j}(u_j^0-v_j^0)^2dx+\sum\limits_j\beta_jd_j\int_0^{l_j}(u_j^3-v_j^3)(u_j-v_j)dx=\\

=[u^0-v^0,u^0-v^0]+\sum\limits_j\beta_jd_j\int_0^{l_j}(u_j-v_j)^2(u_j^2+u_jv_j+v_j^2)dx>0, \text{ если } u^0\neq v^0.

\end{multline*}

Здесь $u=u^0+u^1$, $v=v^0+u^1$, а $[\cdot,\cdot]$ "--- скалярное

произведение, введённое в~рассмотрение при доказательстве леммы

1.1.



Теперь покажем, что оператор $S$ коэрцитивен, т.е.

$$\lim\limits_{||u^0||_\mathfrak{U}\rightarrow\infty}\frac{\langle S(u^0),u^0\rangle}{||u^0||_\mathfrak{U}}=+\infty.$$

Дейст\-ви\-тель\-но,

\begin{multline}\label{svirid:1.5}

\langle S(u^0), u^0\rangle \geq |||

u^0|||^2-|||u^0|||\cdot|||u^1|||+{} \\ +\sum\limits_j\beta_jd_j\int_0^{l_j}\bigl[(u_j^0)^4+3(u_j^0)^3(u_j^1)+

3(u_j^0)^2(u_j^1)^2+u^0(u_j^1)^3\bigr]dx.

\end{multline}

Поскольку все нормы в конечномерном пространстве эквивалентны, то

из (\ref{svirid:1.5}) вытекает коэрцитивность оператора $S$. Здесь через

$|||\cdot|||$ обозначена норма, индуцированная скалярным

произведением $[\cdot, \cdot]$.



Итак, гладкий оператор $S$ строго монотонен и коэрцитивен.

Сле\-до\-ва\-тель\-но, в силу теоремы Вишика"--~Минти"--~Браудера [13, гл. III,

$\S 2$] су\-щест\-ву\-ет единственное решение уравнения $S(u^0)=0$.

Это, в свою очередь, означает, что для любого вектора

$u^1\in{\rm coim \, } L$ существует единственный вектор \mbox{$u^0\in\ker

L$} такой, что $u^0+u^1\in\mathfrak{M}$.

Другими словами, сужение проектора $P$ на $\mathfrak{M}$ "--- биекция.



Теперь пусть точка $u_0\in\mathfrak{M}$. Проверим невырожденность

оператора \linebreak ${(\mathbb{I}-Q)}(M+N'_{u_0})(\mathbb{I}-P)$. Возьмём

вектор $v_0\in\ker L\backslash\{0\}$, тогда

$$\langle(\mathbb{I}-Q)(M+N'_{u_0})v_0,v_0\rangle=|||v_0|||^2+3\sum\limits_j\beta_jd_j\int_0^{l_j}u_0^2v_0^2dx>0.$$



Лемма доказана в случае $\alpha_j\in\mathbb{R}_+$. Если

$\alpha_j\in\mathbb{R}_-$, то для доказательства леммы вместо

оператора $S$ надо взять оператор $T=-S$.

\end{proof}



Из всех рассуждений вытекает

\begin{theorem}[1.1]

Пусть

\begin{itemize}

    \item [\rm (i)] $\ker L=\{0\}$. Тогда фазовым пространством уравнения \eqref{svirid:1.2}

служит все про\-стран\-ство $\mathfrak{U};$



    \item [\rm (ii)] $\ker L\neq\{0\}$, все коэффициенты $\alpha_j\neq\{0\}$ и

имеют тот же знак, что и все ненулевые коэффициенты $\beta_j$. Тогда

фазовым пространством уравнения \eqref{svirid:1.2} служит простое многообразие

$\mathfrak{M}$.

\end{itemize}

\end{theorem}



\begin{remark}[1.1]

Коэффициенты $\lambda_j$ входят в условия

теоремы 1.1 неявным образом, т.\,к. именно они определяют

тривиальность или нетривиальность ядра $\ker L$.

\end{remark}



\begin{remark}[1.2]

Условия на длины $\alpha_j$ и $\beta_j$

хорошо согласуются с условием $\alpha\beta>0$ в [6]. Если

$\alpha\beta<0$, то фазовое пространство уже не будет простым

многообразием [14].

\end{remark}



\Section{Обратная задача}

%\numberwithin{equation}{section}

Для уравнений (\ref{svirid:0.1}) рассмотрим обратную задачу (\ref{svirid:0.2})--(\ref{svirid:0.4}) с

до\-пол\-ни\-тель\-ны\-ми условиями

\begin{gather}\label{svirid:2.1}

\left(\lambda_j+\frac{\partial^2}{\partial x_j^2}\right)u_{jt}(0, 0)=\varphi_j,

\\

\label{svirid:2.2}

\left(\lambda_j+\frac{\partial^2}{\partial x_j^2}\right)u_{jt}(l_j, 0)=\psi_j,

\end{gather}

состоящую в определении неизвестных коэффициентов $\alpha_j$,

$\beta_j$ и решений $u_j$ урав\-не\-ний (\ref{svirid:0.1}).

Касаясь механического смысла задачи, отметим, что здесь $\varphi_j$,

$\psi_j$ показывают изменение скорости динамики выпучивания в~начале

и~конце балки в начальный промежуток времени.



Задачи подобного рода возникают при проведении различных науч\-ных

ис\-сле\-до\-ва\-ний в связи с необходимостью определения характеристик

материала, в котором происходит изучаемый процесс, по результатам

наб\-лю\-де\-ний. Несмотря на различие таких обратных коэффициентных задач, их

объеди\-ня\-ет то, что дополнительная информация в каждой из обратных

за\-дач представляет в общем случае функцию, зависящую от про\-странст\-вен\-ных переменных и

(или) времени.

Такой тип дополнительной

ин\-фор\-ма\-ции характерен для многочисленных экс\-пе\-ри\-мен\-тов в различных

областях нау\-ки и техники.



Подставляя правую часть уравнения (\ref{svirid:0.1}) при $t=0$ в (\ref{svirid:2.1}), (\ref{svirid:2.2}),

получаем системы уравнений относительно $\alpha_j$, $\beta_j$:

\begin{equation}\label{svirid:2.3}

\left \{

\begin{array}{l}

\alpha_j u_{0j}(0)+\beta_j u_{0j}^3(0)=\varphi_j, \\

\alpha_j u_{0j}(l_j)+\beta_ju_{0j}^3(l_j)=\psi_j.

\end{array} \right.

\end{equation}

Для систем уравнений (\ref{svirid:2.3}) воспользуемся правилом Крамера:

\begin{itemize}

    \item [(i)] если

$\mathfrak{\Delta_j}=u_{0j}(0)u_{0j}^3(l_j)-u_{0j}(l_j)u_{0j}^3(0)\neq

0$, то существует единственное ре\-ше\-ние $\alpha_j$, $\beta_j$

каждой из систем уравнений (\ref{svirid:2.3}), причем

\begin{equation}\label{svirid:2.4}

\alpha_j=\frac{\varphi_ju_{0j}^3(l_j)-\psi_ju_{0j}^3(0)}{\mathfrak{\Delta_j}},

\quad

\beta_j=\frac{\psi_ju_{0j}(0)-\varphi_ju_{0j}(l_j)}{\mathfrak{\Delta_j}};

\end{equation}

    \item [(ii)] если $\mathfrak{\Delta_j}=0$, $\varphi_j=\psi_j=0$, то

существует бесконечно много решений $\alpha_j$,~$\beta_j$;

если же $\varphi_j$ или $\psi_j\neq 0$, то решения не существует.

\end{itemize}



Нас интересует случай, когда коэффициенты $\alpha_j$,

$\beta_j\in\mathbb{R}\backslash\{0\}$, поэтому до\-пол\-ни\-тель\-но

предполагаем условия

\begin{equation}\label{svirid:2.5}

\varphi_ju_{0j}^3(l_j)\neq\psi_ju_{0j}^3(0) \quad \mbox{и} \quad

\varphi_ju_{0j}(l_j)\neq\psi_ju_{0j}(0).

\end{equation}

\begin{lemma}[2.1]

Пусть $\ker L=\{0\}.$ Тогда при любых начальных значениях

$u_{0}\in\mathfrak{U},$ $\varphi_j,$ $\psi_j\in\mathbb{R}$ таких{\rm,}

что $u_{0j}(0)\neq 0,$ $u_{0j}(l_j)\neq 0,$ $u_{0j}(0)\neq \pm

u_{0j}(l_j)$ и выполнены условия \eqref{svirid:2.5}, существует единст\-вен\-ное

решение обратной задачи \eqref{svirid:0.1}--\eqref{svirid:0.4}, \eqref{svirid:2.1}--\eqref{svirid:2.2}.

\end{lemma}



\begin{proof}

В силу условий $u_{0j}(0){\neq} 0$, $u_{0j}(l_j){\neq} 0$,

\mbox{$u_{0j}(0){\neq} \pm u_{0j}(l_j)$} определители систем (\ref{svirid:2.3})

$$

\mathfrak{\Delta_j}=u_{0j}(0)u_{0j}^3(l_j)-u_{0j}(l_j)u_{0j}^3(0)=u_{0j}(0)u_{0j}(l_j)(u_{0j}(l_j)-u_{0j}(0))(u_{0j}(l_j)+u_{0j}(0))

$$

не равны нулю, если к тому же выполнены условия (\ref{svirid:2.5}), то коэффициенты $\alpha_j$, $\beta_j\in\mathbb{R}\backslash\{0\}$,

поэтому в~силу теоремы 1.1 решение обратной задачи существует и~единственно для любого вектора $u_{0}\in\mathfrak{U}$.

\end{proof}



Рассмотрим теперь случай нетривиального ядра $\ker L$. Для

выполнения ус\-ло\-вий теоремы 1.1 необходимо, чтобы решения систем

(\ref{svirid:2.3}) были все одного знака. Из (\ref{svirid:2.4}) следует, что для того, чтобы

$\alpha_j$, $\beta_j>0$, должна быть выполнена одна из систем неравенств:



\begin{equation}\label{svirid:2.6}

\left \{

\begin{array}{l}

    \varphi_ju_{0j}^3(l_j)-\psi_ju_{0j}^3(0)>0, \\

    \psi_ju_{0j}(0)-\varphi_ju_{0j}(l_j)>0, \\

    u_{0j}(0)u_{0j}^3(l_j)-u_{0j}(l_j)u_{0j}^3(0)>0

\end{array} \right.

\end{equation}

или

\begin{equation}\label{svirid:2.7}

\left \{

\begin{array}{l}

    \varphi_ju_{0j}^3(l_j)-\psi_ju_{0j}^3(0)<0, \\

    \psi_ju_{0j}(0)-\varphi_ju_{0j}(l_j)<0, \\

    u_{0j}(0)u_{0j}^3(l_j)-u_{0j}(l_j)u_{0j}^3(0)<0,

\end{array} \right.

\end{equation}

а для того, чтобы $\alpha_j$, $\beta_j<0$ "---

\begin{equation}\label{svirid:2.8}

\left \{

\begin{array}{l}

    \varphi_ju_{0j}^3(l_j)-\psi_ju_{0j}^3(0)>0, \\

    \psi_ju_{0j}(0)-\varphi_ju_{0j}(l_j)>0, \\

    u_{0j}(0)u_{0j}^3(l_j)-u_{0j}(l_j)u_{0j}^3(0)<0,

\end{array} \right.

\end{equation}

или

\begin{equation}\label{svirid:2.9}

\left \{

\begin{array}{l}

    \varphi_ju_{0j}^3(l_j)-\psi_ju_{0j}^3(0)<0, \\

    \psi_ju_{0j}(0)-\varphi_ju_{0j}(l_j)<0, \\

    u_{0j}(0)u_{0j}^3(l_j)-u_{0j}(l_j)u_{0j}^3(0)>0.

\end{array} \right.

\end{equation}



Обозначим через $\mathfrak{D}_j^{+}$ область допустимых значений

$\varphi_j$, $\psi_j$, при которых всег\-да верна одна из систем неравенств

(\ref{svirid:2.6}) или (\ref{svirid:2.7}), а через $\mathfrak{D}_j^{-}$ область до\-пус\-ти\-мых

значений $\varphi_j$, $\psi_j$, при которых всегда верна одна из систем

неравенств (\ref{svirid:2.8}) или (\ref{svirid:2.9}). Введём следующие мно\-жест\-ва:

$$ \mathfrak{A}^1_j=\left [

\begin{array}{l}

0<u_{0j}(0)<u_{0j}(l_j), \\

0<-u_{0j}(l_j)<u_{0j}(0), \\

u_{0j}(0)<u_{0j}(l_j)<0, \\

0<-u_{0j}(0)<u_{0j}(l_j),

\end{array} \right. \quad  \mathfrak{A}^2_j=\left [

\begin{array}{l}

0<u_{0j}(l_j)<-u_{0j}(0), \\

u_{0j}(l_j)<u_{0j}(0)<0, \\

0<u_{0j}(l_j)<u_{0j}(0), \\

0<u_{0j}(0)<-u_{0j}(l_j).

\end{array} \right.$$

Легко проверить, что

\begin{equation}\label{svirid:2.10}

\mathfrak{D}_j^{+}= \left \{

\begin{array}{l}

\frac{u_{0j}(l_j)}{u_{0j}(0)}\varphi_j<\psi_j<\left(\frac{u_{0j}(l_j)}{u_{0j}(0)}\right)^3\varphi_j, \mbox{если } u_{0j}(0), u_{0j}(l_j)\in\mathfrak{A}^1_j, \\

\left(\frac{u_{0j}(l_j)}{u_{0j}(0)}\right)^3\varphi_j<\psi_j<\frac{u_{0j}(l_j)}{u_{0j}(0)}\varphi_j,

\mbox{если } u_{0j}(0), u_{0j}(l_j)\in\mathfrak{A}^2_j,

 \end{array} \right.

\end{equation}

\begin{equation}\label{svirid:2.11}

\mathfrak{D}_j^{-}= \left \{

\begin{array}{l}

\left(\frac{u_{0j}(l_j)}{u_{0j}(0)}\right)^3\varphi_j<\psi_j<\frac{u_{0j}(l_j)}{u_{0j}(0)}\varphi_j, \mbox{если } u_{0j}(0), u_{0j}(l_j)\in\mathfrak{A}^1_j, \\

\frac{u_{0j}(l_j)}{u_{0j}(0)}\varphi_j<\psi_j<\left(\frac{u_{0j}(l_j)}{u_{0j}(0)}\right)^3\varphi_j,

\mbox{если } u_{0j}(0), u_{0j}(l_j)\in\mathfrak{A}^2_j.

\end{array} \right.

\end{equation}



Для примера найдем области допустимых значений $\varphi_j$ и $\psi_j$,

при которых коэффициенты $\alpha_j$, $\beta_j>0$ в~случае $\mathfrak{\Delta_j}>0$,

то есть разрешим систему (\ref{svirid:2.6}) при некотором $j$.

В~зависимости от знаков $u_{0j}(0)$ и~$u_{0j}(l_j)$ возможны следующие четыре случая.



\begin{itemize}

    \item [(i)] Пусть $u_{0j}(0)>0$, $u_{0j}(l_j)>0$. Тогда $u_{0j}(0)+u_{0j}(l_j)>0$ и для того, чтобы выполнялось $\mathfrak{\Delta}_j>0$ необходимо $u_{0j}(l_j)-u_{0j}(0)>0$, то есть получился случай $0<u_{0j}(0)<u_{0j}(l_j)$. В этом случае система (\ref{svirid:2.6}) эквивалентна неравенству $$\frac{u_{0j}(l_j)}{u_{0j}(0)}\varphi_j<\psi_j<\frac{u_{0j}^3(l_j)}{u_{0j}^3(0)}\varphi_j.$$



\item [(ii)] Пусть $u_{0j}(0)>0$, $u_{0j}(l_j)<0$. Тогда $u_{0j}(l_j)-u_{0j}(0)<0$ и для того, чтобы выполнялось $\mathfrak{\Delta}_j>0$ необходимо $u_{0j}(l_j)+u_{0j}(0)>0$, то есть получился случай $0<-u_{0j}(l_j)<u_{0j}(0)$. Здесь система (\ref{svirid:2.6}) эквивалентна неравенству $$\frac{u_{0j}(l_j)}{u_{0j}(0)}\varphi_j<\psi_j<\frac{u_{0j}^3(l_j)}{u_{0j}^3(0)}\varphi_j.$$



\item [(iii)] Пусть $u_{0j}(0)<0$, $u_{0j}(l_j)>0$. Тогда $u_{0j}(l_j)-u_{0j}(0)>0$ и для того, чтобы выполнялось $\mathfrak{\Delta}_j>0$ необходимо $u_{0j}(l_j)+u_{0j}(0)<0$, то есть получился случай $0<u_{0j}(l_j)<-u_{0j}(0)$. В этом случае система (\ref{svirid:2.6}) эквивалентна неравенству $$\frac{u_{0j}^3(l_j)}{u_{0j}^3(0)}\varphi_j<\psi_j<\frac{u_{0j}(l_j)}{u_{0j}(0)}\varphi_j.$$

\item [(iv)] Пусть $u_{0j}(0)<0$, $u_{0j}(l_j)<0$. Тогда $u_{0j}(0)+u_{0j}(l_j)<0$ и для того, чтобы выполнялось $\mathfrak{\Delta}_j>0$ необходимо $u_{0j}(l_j)-u_{0j}(0)<0$, то есть получился случай $0<-u_{0j}(0)<-u_{0j}(l_j)$. Здесь система (\ref{svirid:2.6}) эквивалентна неравенству $$\frac{u_{0j}^3(l_j)}{u_{0j}^3(0)}\varphi_j<\psi_j<\frac{u_{0j}(l_j)}{u_{0j}(0)}\varphi_j.$$

\end{itemize}

Аналогично, разрешая системы (\ref{svirid:2.7})--(\ref{svirid:2.9}), получаем (\ref{svirid:2.10}) и (\ref{svirid:2.11}).



Введём множества $\mathfrak{D}^{-}=\left(\mathfrak{D}_1^{-}, \mathfrak{D}_2^{-}, \dots, \mathfrak{D}_j^{-}, \dots\right)$,

$\mathfrak{D}^{+}=\left(\mathfrak{D}_1^{+}, \mathfrak{D}_2^{+}, ..., \mathfrak{D}_j^{+}, ...\right)$ и

$\mathfrak{D}=\mathfrak{D}^{+}\cup\mathfrak{D}^{-}$.

В контексте нашей задачи $\mathfrak{D}^{-}$ "--- множество допустимых значений

векторов $\varphi$, $\psi$, при которых решениями систем (\ref{svirid:2.3})

будут только отрицательные коэффициенты $\alpha_j$, $\beta_j$;

 $\mathfrak{D}^{+}$ "--- множество допустимых значений векторов $\varphi$, $\psi$,

при которых решениями систем (\ref{svirid:2.3}) будут только положительные коэффициенты

$\alpha_j$, $\beta_j$, а $\mathfrak{D}$ "--- множество допустимых значений векторов

$\varphi$, $\psi$, при которых решениями систем (\ref{svirid:2.3}) будут коэффициенты

$\alpha_j$, $\beta_j$ одного знака.



Выберем в ядре $\ker L$ ортонормированный (в смысле

$\langle\cdot,\cdot\rangle$) базис, т.\,е. $\ker L={\rm span \, } \{\chi_k:

k=1, 2, \dots, l\}$, и отождествим его с~базисом в ${\rm coker \, } L$.



Из всего вышесказанного следует

\begin{theorem}[2.1]

\begin{itemize}

  \item [\rm (i)] Пусть $\ker L=\{0\}.$ Тогда при любых

$u_{0}\in\mathfrak{U},$ $\varphi_j,$ $\psi_j\in\mathbb{R}$ таких{\rm ,}

что $u_{0j}(0)\neq 0,$

$u_{0j}(l)\neq 0,$

$u_{0j}(0)\neq \pm u_{0j}(l),$

$\varphi_ju_{0j}^3(l)\neq\psi_ju_{0j}^3(0)$ и

$\varphi_ju_{0j}(l)\neq\psi_ju_{0j}(0),$  существует единственное

решение обратной задачи \eqref{svirid:0.1}--\eqref{svirid:0.4}, \eqref{svirid:2.1}--(\ref{svirid:2.2}).

    \item [\rm(ii)] Пусть $\ker L\neq\{0\}.$

    Тогда при любых начальных данных $\varphi=(\varphi_1, \dots,$ $\varphi_j, \dots)$,

$\psi=(\psi_1, \dots, \psi_j, \dots)\in\mathfrak{D}$

и $u_{0}=(u_{01}, u_{02}, \dots, u_{0j}, \dots)\in\mathfrak{U}$ таких{\rm,}

    что

    \begin{equation}\label{svirid:2.12}

    u_{0j}(0)\neq 0,\quad u_{0j}(l_j)\neq 0,\quad u_{0j}(0)\neq \pm u_{0j}(l_j)

    \end{equation}

    и

\begin{multline}\label{svirid:2.13}

      \left\langle\sum\limits_j\frac{d_j}{\mathfrak{\Delta_j}}\int\limits_0^{l_j}(\varphi_j u_{0j}^3(l_j)-\psi_j u_{0j}^3(0))u_{0j}dx,

    \chi_k\right\rangle+ \\

+\left\langle\sum\limits_j\frac{d_j}{\mathfrak{\Delta_j}}\int\limits_0^{l_j}

(\psi_j u_{0j}(0)-\varphi_j u_{0j}(l_j))u_{0j}^3dx,

\chi_k\right\rangle=0,

\end{multline}

существует единственное решение

$u\in\mathfrak{U},$  $\alpha_j,$

$\beta_j\in\mathbb{R}\backslash\{0\}$,

$\alpha_j\beta_j\in\mathbb{R}_+$ обратной задачи \eqref{svirid:0.1}--\eqref{svirid:0.4},

\eqref{svirid:2.1}--\eqref{svirid:2.2}.

\end{itemize}

\end{theorem}



\begin{proof}

(ii) Пусть $u_0\in\mathfrak{U}$ и выполнены условия (\ref{svirid:2.12}).

Тогда существует единственная пара коэффициентов $\alpha_j$, $\beta_j$,

удовлетворяющая (\ref{svirid:2.3}). Выполнение условий $\varphi$,

$\psi\in\mathfrak{D}$ гарантирует условие $\alpha_j\beta_j>0$ для всех $j$,

причём в~силу конструкции множества $\mathfrak{D}$ все коэффициенты $\alpha_j\neq 0$

и~имеют тот же знак, что и~все ненулевые коэффициенты $\beta_j$.

Тем самым, мы находимся в условиях теоремы 1.1 (ii).

Фазовое пространство

$$

\mathfrak{M}=\{u\in\mathfrak{U}: \left\langle Mu+N(u), \chi_k \right\rangle =0, k=1, 2, \ldots, l\}

$$

при найденных $\alpha_j$ и $\beta_j$ имеет вид (\ref{svirid:2.13}), поэтому существует единственное

ре\-ше\-ние $u\in\mathfrak{U}$,

$\alpha_j$,

$\beta_j\in\mathbb{R}\backslash\{0\}$,

$\alpha_j\beta_j\in\mathbb{R}_{+}$ обратной задачи

(\ref{svirid:0.1})--(\ref{svirid:0.4}), (\ref{svirid:2.1})--(\ref{svirid:2.2}).~\end{proof}



\begin{thebibliography}{99}



\RBibitem{sviridyuk:PYuV}

\by Покорный~Ю.\,В., Пенкин~О.\,М.,  Прядиев~В.\,Л., Боровских~А.\,В., Лазарев~К.\,П., Шабров~С.\,А.

\book Дифференциальные уравнения на геометрических графах

\publaddr М.

\publ Физматлит

\yr 2004

\finalinfo 272~с



\Bibitem{sviridyuk:Kozh}

\by Kozhanov A.\,I.

\book Composite type equations and inverse problems

\publaddr Utrecht

\publ VSP

\yr 1999

\finalinfo 171~pp



\RBibitem{sviridyuk:Svesh}

\book Линейные и нелинейные уравнения соболевского типа

\by Свешников~А.\,Г.,  Альшин~А.\,Б.,  Корпусов~М.\,О., Плетнер~Ю.\,Д.

\publaddr М.

\publ Физматлит

\yr 2007

\finalinfo 736~с





\RBibitem{sviridyuk:Svir}

\by Свиридюк~Г.\,А.

\paper Уравнения соболевского типа на графах

\inbook Неклассические уравнения математической физики

\publaddr Новосибирск

\publ ИМ СО РАН

\yr 2002

\pages 221--225



\Bibitem{sviridyuk:Fed}

\by Fedorov~V.\,E., Urazaeva~A.\,V.

\paper  An inverse problem for linear Sobolev type equations

\jour J.~Inv.~Ill-Posed Problems

\yr 2004

\vol 12

\issue 5

\pages 1--9



\RBibitem{sviridyuk:SvB}

\by Свиридюк~Г.\,А., Баязитова~А.\,А.

\paper Обратная задача для уравнений Ба\-рен\-блат\-та"--~Желтова"--~Кочиной на графе

\inbook Неклассические уравнения математической физики

\bookinfo Тр. междунар. конф. <<Дифференциальные уравнения, теория функ\-ций и приложения>>, посвящ. 100-летию со дня рождения

акад. И.\,Н.~Векуа

\publaddr Новосибирск

\publ  ИМ СО РАН

\yr 2007

\pages 244--250



\RBibitem{sviridyuk:SvSh}

\paper Уравнения Хоффа на графах

\by Свиридюк~Г.\,А., Шеметова~В.\,В.

\jour Дифференц. уравнения

\yr 2006

\vol 42

\issue 1

\pages 126--131





\Bibitem{sviridyuk:SvF}

\book Sobolev type equations and degenerate semi\-groups of operators

\by  Sviridyuk~G.\,A., Fedorov~V.\,E.

\publaddr Utrecht; Boston

\publ VSP

\yr 2003

\finalinfo 216~pp





\RBibitem{sviridyuk:L}

\by Ленг~С.

\book Введение в теорию дифференцируемых многообразий

\publaddr Волгоград

\publ Платон

\yr 1996

\finalinfo 203~с



\RBibitem{sviridyuk:Svi}

\by Свиридюк~Г.\,А.

\paper Многообразие решений одного нелинейного сингулярного псевдопараболического уравнения

\jour Докл. АН СССР

\yr 1986

\vol 289

\issue 6

\pages 1315--1318



\RBibitem{sviridyuk:Sv}

\by Свиридюк, Г.\,А.

\paper Одна задача для обобщенного фильтрационного уравнения Буссинеска

\jour Изв. вузов. Математика

\yr 1989

\issue 2

\pages 55--61



\RBibitem{sviridyuk:SvM}

\paper Задача оптимального управления для уравнения Хоф\-фа

\by Свиридюк~Г.\,А., Манакова~Н.\,А.

\jour  Сиб. журн. индустр. математики

\yr 2005

\vol 8

\issue 2

\pages 144--151





\RBibitem{sviridyuk:GGZ}

\book Нелинейные операторные уравнения и опе\-ра\-тор\-ные дифференциальные уравнения

\by Гаевский~Х., Грегер~К., Захариас~К.

\publaddr М.

\publ Мир

\yr 1978

\finalinfo 336~c



\RBibitem{sviridyuk:SvT}

\paper Сборка Уитни в фазовом пространстве уравнения Хоффа

\by Свиридюк~Г.\,А., Тринеева~И.\,К.

\jour Изв. вузов. Математика

\yr 2005

\issue 10

\pages 54--60



\end{thebibliography}





\makeengtitlem







